\documentclass[aps,prd,reprint,nofootinbib]{revtex4-2}
% Co-theoretical-physicist paper skeleton.
% Compile with pdflatex (revtex4-2). Swap prd -> prl/prb/pre as appropriate.

\usepackage{amsmath,amssymb,bm}
\usepackage{graphicx}

\begin{document}

\title{<Title>}
\author{<author>}
\affiliation{<affiliation>}
\date{\today}

\begin{abstract}
<Problem, the core idea/unifying principle, main results, and how they were
established (analytic + numerical), in one paragraph.>
\end{abstract}

\maketitle

\section{Introduction}
% The problem, why it matters, what was known (cite physics-literature), and the
% contribution of this work.

\section{Framework}
% Effective description and degrees of freedom; assumptions; symmetries; the model.

\section{Results}
% Derived results (from physics-derivation) and numerically-established results
% (from physics-numerics). For each, state its evidence-ladder status.

\section{Predictions}
% New falsifiable predictions and the experiments/numerics that would test them.

\section{Domain of validity}
% Examples where the theory applies; non-examples / where it breaks.

\section{Discussion}
% Relation to prior work, open questions, next steps.

% --- Dependency map (move to a comment or appendix) ---
% result R1: depends on {assumptions A1}, status = analytically-derived
% result R2: depends on {R1, numerics N3},  status = numerically-supported
% ...

\bibliography{refs}

\end{document}
